\section{Conclusion}\label{sec:conclusion}

We asked where jagged competence comes from. In a task built from two known
layers, the answer was that a network trained on the upper layer learns the
lower one only in a jagged form, readable where the upper layer needs it,
with B's category responses failing near the cutoffs, and not reliably
delivered or kept; these separations explain part of its failures, and some
remain unexplained. Knowing an intermediate quantity
turned out to be several separate achievements, each of which can fail on its
own, and average error and probes each hid the failures.

The practical lesson is that a network's competence on a task built in
layers should not be judged by its average error or by what a probe can read,
but by measuring, layer by layer, what it has acquired, delivered and kept.
Whether learning an upper theory yields a jagged lower theory in larger
systems is a test to run.
